% ========================================================
% BÀI TẬP TOÁN 5 - KẾT NỐI TRI THỨC VỚI CUỘC SỐNG
% Bài 10: Khái niệm số thập phân
% Định dạng: BTTeX Part III - Trắc nghiệm điền số ngắn
% ========================================================

% Câu 1
\begin{ex}
  Trong số thập phân $35{,}74$, chữ số ở hàng phần mười là bao nhiêu?
  \shortans{7}
  \loigiai{
    Trong số $35{,}74$, chữ số ngay sau dấu phẩy là chữ số 7 thuộc hàng phần mười.
  }
\end{ex}

% Câu 2
\begin{ex}
  Trong số thập phân $128{,}065$, chữ số ở hàng phần trăm là bao nhiêu?
  \shortans{6}
  \loigiai{
    Chữ số ở hàng phần mười là 0, hàng phần trăm là 6.
  }
\end{ex}

% Câu 3
\begin{ex}
  Trong số thập phân $49{,}382$, chữ số ở hàng phần nghìn là bao nhiêu?
  \shortans{2}
  \loigiai{
    Chữ số 2 đứng ở vị trí thứ ba sau dấu phẩy, thuộc hàng phần nghìn.
  }
\end{ex}

% Câu 4
\begin{ex}
  Trong số thập phân $872{,}15$, chữ số $8$ có giá trị bằng bao nhiêu?
  \shortans{800}
  \loigiai{
    Chữ số 8 thuộc hàng trăm nên có giá trị là 800.
  }
\end{ex}

% Câu 5
\begin{ex}
  Phần nguyên của số thập phân $205{,}48$ là bao nhiêu?
  \shortans{205}
  \loigiai{
    Phần nguyên là phần đứng trước dấu phẩy, bằng 205.
  }
\end{ex}

% Câu 6
\begin{ex}
  Chuyển phân số thập phân $\frac{7}{10}$ thành số thập phân. Phần thập phân của số đó có giá trị là bao nhiêu phần mười?
  \shortans{7}
  \loigiai{
    $\frac{7}{10} = 0{,}7$. Phần thập phân gồm 7 phần mười.
  }
\end{ex}

% Câu 7
\begin{ex}
  Chuyển phân số thập phân $\frac{45}{100}$ thành số thập phân $0{,}45$. Chữ số hàng phần trăm là bao nhiêu?
  \shortans{5}
  \loigiai{
    Chữ số hàng phần trăm là 5.
  }
\end{ex}

% Câu 8
\begin{ex}
  Chuyển phân số thập phân $\frac{125}{1000}$ thành số thập phân $0{,}125$. Tổng các chữ số ở phần thập phân là bao nhiêu?
  \shortans{8}
  \loigiai{
    Các chữ số ở phần thập phân là 1, 2, 5. Tổng là $1 + 2 + 5 = 8$.
  }
\end{ex}

% Câu 9
\begin{ex}
  Viết số thập phân gồm $5$ đơn vị và $6$ phần mười dưới dạng $5{,}6$. Tích của phần nguyên và chữ số phần mười là bao nhiêu?
  \shortans{30}
  \loigiai{
    Số đó là $5{,}6$. Tích là $5 \times 6 = 30$.
  }
\end{ex}

% Câu 10
\begin{ex}
  Viết số thập phân gồm $12$ đơn vị, $3$ phần mười và $8$ phần trăm. Tổng các chữ số của số thập phân đó là bao nhiêu?
  \shortans{14}
  \loigiai{
    Số thập phân đó là $12{,}38$. Tổng các chữ số là $1 + 2 + 3 + 8 = 14$.
  }
\end{ex}

% Câu 11
\begin{ex}
  Số thập phân $0{,}07$ có bao nhiêu chữ số ở phần thập phân?
  \shortans{2}
  \loigiai{
    Phần thập phân là $07$ có 2 chữ số.
  }
\end{ex}

% Câu 12
\begin{ex}
  Số thập phân gồm $8$ chục, $5$ đơn vị và $9$ phần nghìn được viết là $85{,}009$. Chữ số hàng phần trăm của số này là bao nhiêu?
  \shortans{0}
  \loigiai{
    Số là $85{,}009$. Chữ số hàng phần trăm là 0.
  }
\end{ex}

% Câu 13
\begin{ex}
  Khi viết thêm $2$ chữ số $0$ vào bên phải phần thập phân của số $4{,}5$ ta được số $4{,}500$. Giá trị của số đó thay đổi bao nhiêu đơn vị?
  \shortans{0}
  \loigiai{
    Khi viết thêm chữ số 0 vào tận cùng bên phải phần thập phân thì giá trị của số không thay đổi, tức là thay đổi 0 đơn vị.
  }
\end{ex}

% Câu 14
\begin{ex}
  Có bao nhiêu chữ số $0$ có thể xóa đi ở tận cùng bên phải phần thập phân của số $12{,}3000$ mà giá trị không đổi?
  \shortans{3}
  \loigiai{
    Có thể xóa đi 3 chữ số 0 ở cuối để được số $12{,}3$.
  }
\end{ex}

% Câu 15
\begin{ex}
  Chuyển hỗn số $3\frac{7}{10}$ thành số thập phân ta được $3{,}7$. Tổng của phần nguyên và chữ số hàng phần mười là bao nhiêu?
  \shortans{10}
  \loigiai{
    $3\frac{7}{10} = 3{,}7$. Tổng là $3 + 7 = 10$.
  }
\end{ex}

% Câu 16
\begin{ex}
  Chuyển hỗn số $15\frac{24}{100}$ thành số thập phân. Chữ số hàng phần mười của số này là bao nhiêu?
  \shortans{2}
  \loigiai{
    $15\frac{24}{100} = 15{,}24$. Chữ số hàng phần mười là 2.
  }
\end{ex}

% Câu 17
\begin{ex}
  Viết phân số $\frac{3}{4}$ thành số thập phân ta được $0{,}75$. Tổng chữ số hàng phần mười và hàng phần trăm là bao nhiêu?
  \shortans{12}
  \loigiai{
    $\frac{3}{4} = \frac{75}{100} = 0{,}75$. Tổng là $7 + 5 = 12$.
  }
\end{ex}

% Câu 18
\begin{ex}
  Viết phân số $\frac{1}{2}$ thành số thập phân ta được $0{,}5$. Lấy phần nguyên cộng với chữ số hàng phần mười được bao nhiêu?
  \shortans{5}
  \loigiai{
    Phần nguyên là 0, hàng phần mười là 5. Tổng là $0 + 5 = 5$.
  }
\end{ex}

% Câu 19
\begin{ex}
  Trong số thập phân $67{,}89$, hiệu giữa chữ số ở hàng phần trăm và chữ số ở hàng phần mười là bao nhiêu?
  \shortans{1}
  \loigiai{
    Chữ số hàng phần trăm là 9, hàng phần mười là 8. Hiệu là $9 - 8 = 1$.
  }
\end{ex}

% Câu 20
\begin{ex}
  Số thập phân nào có phần nguyên là số lẻ bé nhất có hai chữ số, phần thập phân là số chẵn lớn nhất có một chữ số? (Viết tổng các chữ số của số đó).
  \shortans{10}
  \loigiai{
    Phần nguyên là 11. Phần thập phân là 8. Số đó là $11{,}8$. Tổng các chữ số là $1 + 1 + 8 = 10$.
  }
\end{ex}

% Câu 21
\begin{ex}
  Tìm chữ số $x$ trong số thập phân $2{,}x4$ biết chữ số hàng phần mười $x$ gấp đôi chữ số hàng phần trăm.
  \shortans{8}
  \loigiai{
    Chữ số hàng phần trăm là 4. Vì chữ số hàng phần mười gấp đôi chữ số hàng phần trăm nên $x = 4 \times 2 = 8$.
  }
\end{ex}

% Câu 22
\begin{ex}
  Chuyển số thập phân $0{,}6$ thành phân số tối giản $\frac{a}{b}$. Mẫu số $b$ là bao nhiêu?
  \shortans{5}
  \loigiai{
    $0{,}6 = \frac{6}{10} = \frac{3}{5}$. Mẫu số là 5.
  }
\end{ex}

% Câu 23
\begin{ex}
  Chuyển số thập phân $0{,}25$ thành phân số tối giản $\frac{a}{b}$. Tổng tử số và mẫu số $a + b$ bằng bao nhiêu?
  \shortans{5}
  \loigiai{
    $0{,}25 = \frac{25}{100} = \frac{1}{4}$. Tổng là $1 + 4 = 5$.
  }
\end{ex}

% Câu 24
\begin{ex}
  Một thước dây dài $1\text{ m } 8\text{ dm}$. Viết độ dài thước dây dưới dạng số thập phân có đơn vị mét là $1{,}8\text{ m}$. Chữ số hàng phần mười là bao nhiêu?
  \shortans{8}
  \loigiai{
    Độ dài là $1{,}8\text{ m}$. Chữ số hàng phần mười là 8.
  }
\end{ex}

% Câu 25
\begin{ex}
  Một bao gạo nặng $45\text{ kg } 500\text{ g}$. Viết khối lượng dưới dạng số thập phân với đơn vị ki-lô-gam là $45{,}5\text{ kg}$. Chữ số ở hàng đơn vị là bao nhiêu?
  \shortans{5}
  \loigiai{
    Khối lượng là $45{,}5\text{ kg}$. Chữ số hàng đơn vị là 5.
  }
\end{ex}
